% ============================================================
% CAPITOLO 7 - ERRORI DI CONDOTTA DELLE VOCI
% ============================================================
\chapter{Errori di condotta delle voci}
\label{chap:errori_condotta}

Gli errori di condotta si dividono in melodici, che riguardano una voce presa da sola, e armonici, che nascono dal rapporto tra due voci in successione. Questo capitolo li elenca e dice quali eccezioni la pratica ammette.

\section{Errori melodici e falsa relazione}
\label{sec:errori_melodici}

Ogni linea vocale deve restare cantabile. Per questo la condotta di una singola voce ha dei limiti.

\begin{regola}[Limiti degli intervalli melodici]
	\label{reg:limiti_melodici}
	Nella condotta di una singola voce valgono queste norme:
	\begin{enumerate}
		\item Sono \textbf{proibiti} gli intervalli superiori alla \textbf{sesta minore}, eccetto l'ottava giusta.
		\item Sono \textbf{proibiti} tutti gli intervalli \textbf{eccedenti e diminuiti}, con le eccezioni indicate sotto.
		\item La \textbf{settima minore ascendente} è ammessa in due casi: quando sostituisce il movimento di tono discendente tra il V e il IV grado, e in funzione modulante.
	\end{enumerate}
\end{regola}

\begin{osservazione}[Eccezione: seconda aumentata nel minore]
	\label{oss:seconda_aumentata}
	La \textbf{seconda aumentata} tra il VI e il VII grado della scala minore armonica è \textit{tollerata} quando non si può evitare. Nasce dalla presenza simultanea della sopradominante naturale e della sensibile alzata.
\end{osservazione}

\begin{osservazione}[Eccezione: quarta e quinta diminuita]
	\label{oss:quarta_quinta_dim}
	Gli intervalli di \textbf{quarta e quinta diminuita} sono ammessi se il movimento è \textbf{discendente} e la nota successiva risale di \textbf{semitono diatonico}. È lo schema tipico del VII grado che funge da appoggiatura.
\end{osservazione}

\begin{errore}[Quarta eccedente melodica]
	\label{err:quarta_eccedente}
	Una voce non può muoversi per salto di \textbf{quarta eccedente} (tritono). L'intervallo, chiamato storicamente \textit{diabolus in musica}, è difficile da cantare e spezza la linea melodica.
\end{errore}

\begin{regola}[Eccezione: quarta eccedente con risoluzione]
	\label{reg:quarta_eccedente_eccezione}
	Il salto di quarta eccedente è ammesso se la nota d'arrivo prosegue di \textbf{semitono diatonico} nella \textbf{stessa direzione} del salto (Fa-Si-Do salendo, Si-Fa-Mi scendendo): il tritono risolve così verso l'esterno, come fa nell'accordo.
\end{regola}

\subsection{La falsa relazione}
\label{subsec:falsa_relazione}

La falsa relazione è un errore di natura orizzontale che si manifesta fra due voci \textit{diverse}.

\begin{definizione}[Falsa relazione]
	\label{def:falsa_relazione}
	Si ha \textbf{falsa relazione} quando la stessa nota compare alterata in una voce e non alterata in un'altra, in accordi consecutivi. Distrugge la coerenza cromatica del tessuto.
	\begin{itemize}
		\item \textbf{Falsa relazione melodica:} un suono viene alterato in una voce diversa da quella che l'ha cantato non alterato nell'accordo precedente.
		\item \textbf{Falsa relazione di tritono:} le due note del tritono diatonico (IV e VII grado) compaiono in due voci diverse in accordi consecutivi, creando una relazione ambigua tra un accordo e il successivo.
	\end{itemize}
\end{definizione}

\begin{regola}[Come evitare la falsa relazione]
	\label{reg:evitare_falsa_relazione}
	Se una nota viene alterata, la stessa voce che l'ha cantata non alterata deve cantare la versione alterata nell'accordo successivo. L'alterazione non si «passa» a un'altra voce.
\end{regola}

\section{Quinte e ottave parallele}
\label{sec:quinte_ottave_parallele}

Gli errori armonici gravi sono i parallelismi di quinte e di ottave.

\begin{errore}[Quinte e ottave parallele]
	\label{err:quinte_ottave_parallele}
	È vietato muovere due voci per moto parallelo mantenendo tra loro un intervallo di \textbf{quinta giusta} o di \textbf{ottava giusta}. Il parallelismo toglie alle parti l'indipendenza, e le fa percepire come un'unica linea raddoppiata.
	\musica{es_cap04_03}
\end{errore}

Due quinte o due ottave sono ammesse se le separa un'armonia diversa e non seguono quindi in successione diretta.

\begin{regola}[Quinte parallele giusta e diminuita]
	\label{reg:quinte_g_dim}
	Il moto di quinte parallele è ammesso quando la \textbf{prima quinta è giusta} e la \textbf{seconda è diminuita}. La quinta diminuita, essendo dissonante, non è percepita come una vera quinta parallela.
	\musica{es_cap04_04}
\end{regola}

\section{Quinte e ottave nascoste}
\label{sec:quinte_nascoste}

Gli errori armonici meno gravi sono i parallelismi nascosti, o diretti.

\begin{errore}[Quinte e ottave nascoste (o dirette)]
	\label{err:quinte_nascoste}
	Si deve evitare di arrivare a un'ottava o a una quinta giusta tra le parti estreme (soprano e basso) muovendo entrambe le voci nella stessa direzione per salto.
\end{errore}

\begin{regola}[Eccezioni alle quinte e ottave nascoste]
	\label{reg:eccezioni_nascoste}
	Le quinte e ottave nascoste tra le parti estreme sono tollerate in questi casi:
	\begin{itemize}
		\item per le \textbf{ottave nascoste}: il basso sale di almeno una quarta giusta e il soprano sale di semitono diatonico (tipicamente la sensibile che sale alla tonica);
		\item per le \textbf{quinte nascoste}: una delle voci procede per grado congiunto, oppure il basso va verso la tonica o la dominante e la voce superiore procede per grado.
	\end{itemize}
\end{regola}

\section{Incrocio delle parti}
\label{sec:incrocio_parti}

\begin{errore}[Incrocio delle voci]
	\label{err:incrocio}
	Una voce più grave non può cantare una nota più acuta di quella della voce adiacente superiore: il tenore non supera il contralto, né il contralto il soprano.
\end{errore}

\section{La sincope armonica}
\label{sec:sincope_armonica}

La \textbf{sincope armonica} è un caso particolare di condotta strettamente legato ai ritardi (Capitolo \ref{chap:ritardi}): una voce tiene una nota che nell'accordo successivo diventa dissonante.

\begin{definizione}[Sincope armonica]
	\label{def:sincope_armonica}
	Si ha \textbf{sincope armonica} quando una nota \textit{consonante} rispetto all'accordo corrente viene \textit{tenuta} mentre il basso si muove verso un nuovo accordo, e diventa così \textit{dissonante} rispetto alla nuova armonia.
\end{definizione}

\begin{regola}[Trattamento della sincope armonica]
	\label{reg:sincope_armonica}
	La nota tenuta per sincope deve:
	\begin{enumerate}
		\item essere preparata come consonanza nell'accordo precedente;
		\item cadere sul tempo forte, o su una parte forte del tempo, che rafforza la tensione;
		\item risolvere per grado congiunto \textbf{discendente} su una nota reale dell'accordo, comportandosi come un ritardo.
	\end{enumerate}
\end{regola}
